\documentclass[12pt]{article}
\usepackage[spanish]{babel}
\usepackage[utf8]{inputenc}
\usepackage{amsmath}
\usepackage{amssymb}
\usepackage{array}
\usepackage{geometry}
\usepackage{graphicx}
\usepackage{booktabs}

\geometry{a4paper, margin=2.5cm}

\title{Práctica: Interpolación de Newton y Diferencias Finitas}
\author{Métodos Numéricos}
\date{}

\begin{document}
\maketitle

\section*{Descripción del problema}

Una persona de altura $1.86 \, \text{m}$ lanza una pelota en tres ocasiones distintas. 
Se registra la posición vertical de la pelota (en metros) en función del tiempo (en segundos). 
Estos datos representan la posición de la partícula en función del tiempo.

\section*{Datos experimentales}

\subsection*{Lanzamiento 1}
\begin{center}
\begin{tabular}{cc}
\toprule
$t$ (s) & $y$ (m) \\
\midrule
0.0 & 1.86 \\
0.5 & 3.20 \\
1.0 & 4.10 \\
1.5 & 4.45 \\
2.0 & 4.20 \\
2.5 & 3.40 \\
3.0 & 2.10 \\
\bottomrule
\end{tabular}
\end{center}

\subsection*{Lanzamiento 2}
\begin{center}
\begin{tabular}{cc}
\toprule
$t$ (s) & $y$ (m) \\
\midrule
0.0 & 1.86 \\
0.5 & 3.00 \\
1.0 & 3.95 \\
1.5 & 4.60 \\
2.0 & 4.55 \\
2.5 & 3.70 \\
3.0 & 2.40 \\
\bottomrule
\end{tabular}
\end{center}

\subsection*{Lanzamiento 3}
\begin{center}
\begin{tabular}{cc}
\toprule
$t$ (s) & $y$ (m) \\
\midrule
0.0 & 1.86 \\
0.5 & 2.80 \\
1.0 & 3.60 \\
1.5 & 4.05 \\
2.0 & 3.90 \\
2.5 & 3.10 \\
3.0 & 1.95 \\
\bottomrule
\end{tabular}
\end{center}

\section*{Actividades}

\begin{enumerate}
\item \textbf{Interpolación de Newton}
\begin{itemize}
\item Para cada lanzamiento, construya el polinomio interpolante de Newton usando todos los puntos.
\item Grafique los datos y el polinomio interpolante.
%\item Interprete físicamente el significado del polinomio obtenido.
\end{itemize}

\item \textbf{Velocidad mediante diferencias finitas}

Usando los datos del lanzamiento 1, calcule la velocidad en $t = 1.5$ s utilizando:

\begin{itemize}
\item Diferencias finitas de orden 4
\item Diferencias finitas de orden 5
\item Diferencias finitas de orden 6
\end{itemize}

Compare los resultados obtenidos.

\item \textbf{Análisis del error}
\begin{itemize}
\item Explique cómo cambia el error al aumentar el orden del método.
\item ¿Siempre aumentar el orden mejora la aproximación? Cuidado : Aunque un método tenga mayor orden, no siempre produce menor error, porque los métodos de orden alto amplifican el ruido de los datos.
\end{itemize}

\item \textbf{Interpretación física}
\begin{itemize}
\item ¿En qué instante la velocidad es cero?
\item ¿Qué representa ese punto físicamente?
\item Compare los tres lanzamientos e indique cuál tuvo mayor velocidad inicial.
\item Es posible que algunos de los lanzamiento alcance a un cuerpo ubicado a 3 mts del lanzador que tiene una altura de 1.50 mts
\end{itemize}
\end{enumerate}

\section*{Tabla de diferencias finitas}

Sea una función $y(t)$ con paso uniforme $h$. La tabla de diferencias finitas se construye de la siguiente manera:

\begin{center}
\begin{tabular}{c|c|c|c|c|c}
$t$ & $y$ & $\Delta y$ & $\Delta^2 y$ & $\Delta^3 y$ & $\Delta^4 y$ \\
\hline
$t_0$ & $y_0$ & $\Delta y_0$ & $\Delta^2 y_0$ & $\Delta^3 y_0$ & $\Delta^4 y_0$ \\
$t_1$ & $y_1$ & $\Delta y_1$ & $\Delta^2 y_1$ & $\Delta^3 y_1$ &  \\
$t_2$ & $y_2$ & $\Delta y_2$ & $\Delta^2 y_2$ &  &  \\
$t_3$ & $y_3$ & $\Delta y_3$ &  &  &  \\
$t_4$ & $y_4$ &  &  &  &  \\
$t_5$ & $y_5$ &  &  &  &  \\
\end{tabular}
\end{center}

Donde:

\[
\Delta y_i = y_{i+1} - y_i
\]

\[
\Delta^2 y_i = \Delta y_{i+1} - \Delta y_i
\]

\[
\Delta^3 y_i = \Delta^2 y_{i+1} - \Delta^2 y_i
\]

\[
\Delta^4 y_i = \Delta^3 y_{i+1} - \Delta^3 y_i
\]


\section*{Fórmulas de derivación numérica}

Sea $h$ el paso entre los datos.

\subsection*{Derivada centrada orden 4}

\[
f'(x_i) \approx 
\frac{-f(x_{i+2}) + 8f(x_{i+1}) - 8f(x_{i-1}) + f(x_{i-2})}{12h}
\]

Error:

\[
E = O(h^4)
\]


\subsection*{Derivada centrada orden 5}

\[
f'(x_i) \approx
\frac{-2f(x_{i+3}) + 15f(x_{i+2}) - 60f(x_{i+1}) + 20f(x_i) + 30f(x_{i-1}) - 3f(x_{i-2})}{60h}
\]

Error:

\[
E = O(h^5)
\]


\subsection*{Derivada centrada orden 6}

\[
f'(x_i) \approx
\frac{f(x_{i-3}) - 9f(x_{i-2}) + 45f(x_{i-1}) - 45f(x_{i+1}) + 9f(x_{i+2}) - f(x_{i+3})}{60h}
\]

Error:

\[
E = O(h^6)
\]



\end{document}